\chapter{Αρχιτεκτονική λογισμικού}

\section{Αρχιτεκτονική και δομή ETL}

Η διαδικασία ETL υλοποιείται εξ ολοκλήρου σε Python με τη βιβλιοθήκη
\emph{Polars}. Δεν χρησιμοποιεί ORM ούτε αγγίζει τη βάση: τερματίζει με την
παραγωγή CSV αρχείων στο φάκελο \texttt{exports/}. Η ενοποίηση γίνεται σε
τέσσερις διακριτές φάσεις.

\paragraph{(α) Extract.} Τα πηγαία CSV αρχεία (DBLP journals \& conferences,
iCore26, Kaggle Scimago) διαβάζονται μέσα σε
\texttt{polars.LazyFrame} για streaming επεξεργασία.

\paragraph{(β) Transform.}

\begin{itemize}
    \item \textbf{Normalisation}: lowercase, expansion συντομογραφιών
          (\texttt{Trans. -> Transactions}, \texttt{Comput. -> Computing}),
          αφαίρεση τόνων.
    \item \textbf{Entity resolution} (το πιο κρίσιμο βήμα):
          \begin{itemize}
              \item \textbf{Conferences}: case-insensitive exact join του
                    DBLP \texttt{booktitle} (που συνήθως είναι το acronym
                    \texttt{EDBT, ICDE, VLDB \dots}) με το iCore26 \texttt{Acronym}.
                    Recall \(\approx\) 19\%, precision 100\%.
              \item \textbf{Journals}: \texttt{rapidfuzz.token\_set\_ratio} με
                    threshold 85 πάνω σε pre-normalised titles. Ετσι το
                    DBLP \texttt{IEEE Trans. Knowl. Data Eng.} ενώνεται με
                    το Kaggle \texttt{IEEE Transactions on Knowledge and
                    Data Engineering}. Recall \(\approx\) 80\%.
              \item \textbf{Authors}: \texttt{|}-split του DBLP authors
                    field, αφαίρεση leading/trailing whitespace, exact
                    matching (έλλειψη δεδομένων για συνωνυμίες).
          \end{itemize}
    \item \textbf{Validation}: γραμμές με NULL title / year εκτός εύρους
          [1900, 2030] / άδειο authors field απορρίπτονται και γράφονται
          στο \texttt{rejection\_log.csv} με reason και raw row JSON
          ώστε να μπορούν να ανακατασκευαστούν αν ξανατρέξει το ETL με
          βελτιωμένη λογική.
    \item \textbf{Surrogate keys}: σε όλους τους lookup πίνακες
          (\texttt{author\_id}, \texttt{journal\_id}, \texttt{conference\_id})
          για να συμμορφωθούμε με το \emph{integer-PK} requirement.
\end{itemize}

\paragraph{(γ) Load.} Τα CSVs γράφονται με
\texttt{Polars.write\_csv(null\_value="\textbackslash N")} ώστε το \texttt{\textbackslash N}
sentinel να μεταφραστεί σε \texttt{NULL} από τον \texttt{LOAD DATA INFILE} χωρίς
extra preprocessing. Το \texttt{sql\_scripts/02\_load.sql} εκτελεί τη φόρτωση σε
σωστή σειρά (lookups πρώτα, fact tables μετά, bridges τελευταία).

\paragraph{(δ) Post-load.} Το \texttt{sql\_scripts/03\_views.sql} εκτελείται μετά
το load και δημιουργεί τα views και τα materialised tables (Section
\ref{sec:physical}).

\paragraph{Reproducible quality audit.}
Όλα τα ποσοστά καθαρισμού που αναφέρονται στην παρούσα αναφορά
παράγονται απευθείας από το script
\texttt{scripts/data\_quality\_report.sql}. Το script είναι
αναπαραγωγίσιμο σε οποιοδήποτε στιγμιότυπο της βάσης και δίνει:
row counts ανά πίνακα, ποσοστά match (journals
$\approx 80.2$\%, conferences $\approx 19.0$\%), breakdown
απορρίψεων από το \texttt{rejection\_logs}, και integrity checks
(orphan articles / bridges) που πρέπει όλα να επιστρέφουν 0 μετά
από καθαρή φόρτωση. Εκτελείται ως:

\begin{lstlisting}[language=bash]
docker exec -i mye030_mysql mysql -uroot -proot mye030 \
    < scripts/data_quality_report.sql
\end{lstlisting}

Έτσι κάθε ποσοστό είναι ελέγξιμο αντί να βασίζεται σε χειρόγραφο
ad-hoc query.

\begin{figure}[H]
    \centering
    \includegraphics[width=\linewidth,height=0.78\textheight,keepaspectratio]{etl_activity.png}
    \caption{Component / activity diagram του ETL pipeline. Τα DBLP, iCore26,
    Kaggle CSV εισέρχονται στον \texttt{transformer.py}, περνούν από
    normalisation, entity resolution, validation, βγαίνουν σε
    \texttt{exports/}\,*.csv και rejection\_log.csv, φορτώνονται στη MySQL
    μέσω \texttt{02\_load.sql} και τέλος το \texttt{03\_views.sql} δημιουργεί
    views και materialised tables.}
    \label{fig:etl}
\end{figure}

\section{Διαγράμματα πακέτων / υποσυστημάτων κεντρικής εφαρμογής}

\subsection{Back-end packages}

\begin{itemize}
    \item \texttt{api.routers}: HTTP layer (FastAPI). Ενα module ανά resource
          (\texttt{journals.py}, \texttt{conferences.py}, \texttt{authors.py},
          \texttt{years.py}, \texttt{charts.py}, \texttt{meta.py}). 28 endpoints
          συνολικά (5 + 5 + 4 + 5 + 7 + 2).
    \item \texttt{api.data\_access}: ένα DAO ανά resource με raw, parameterised
          SQL queries. Κάθε method χρησιμοποιεί \texttt{Database.fetch\_all} /
          \texttt{fetch\_one} / \texttt{execute}.
    \item \texttt{api.schemas}: Pydantic v2 BaseModels για όλες τις request /
          response DTOs. Δεν χρησιμοποιούνται ως ORM models.
    \item \texttt{api.database}: λεπτό async wrapper γύρω από
          \texttt{asyncmy} pool.
    \item \texttt{api.dependencies}, \texttt{api.settings}: dependency injection
          και environment-driven configuration.
\end{itemize}

\begin{figure}[H]
    \centering
    \includegraphics[width=0.85\linewidth,height=0.78\textheight,keepaspectratio]{packages_backend.png}
    \caption{Package diagram του back-end: routers $\to$ data\_access $\to$
    database $\to$ MySQL. Καμία direct εξάρτηση από τα DAOs προς routers
    και ξεκάθαρη μονόπλευρη ροή.}
    \label{fig:packages-backend}
\end{figure}

\subsection{Front-end packages}

\begin{itemize}
    \item \texttt{src/routes/}: 13 file-based routes (TanStack Router) σε
          ένα προς ένα αντιστοιχία με τις σελίδες της εφαρμογής (Landing,
          Dashboard, Journals/Conferences/Authors/Years
          list+profile, Paper details, Charts playground).
    \item \texttt{src/components/charts/}: 6 Visx-based chart components
          (\texttt{LineChart}, \texttt{BarChart}, \texttt{ScatterPlot},
          \texttt{Heatmap}, \texttt{StackedArea}, \texttt{HorizontalBarChart})
          με κοινό \texttt{ChartFrame} wrapper και shared \texttt{InteractiveLegend}.
    \item \texttt{src/components/filters/}: \texttt{SearchInput} (debounced),
          \texttt{SelectFilter}, \texttt{YearRangeFilter},
          \texttt{CheckboxList} (scrolling multi-select),
          \texttt{VenuePicker}.
    \item \texttt{src/components/tables/}: \texttt{PaginatedTable} με sortable
          columns, sticky header, zebra striping.
    \item \texttt{src/api/}: TanStack Query hooks γύρω από
          \texttt{apiFetch}· τύποι DTOs αντιγραμμένοι χειροκίνητα από τα
          Pydantic schemas (single source of truth: το OpenAPI).
\end{itemize}

\section{Deployment diagram}

\begin{figure}[H]
    \centering
    \includegraphics[width=0.95\linewidth,height=0.78\textheight,keepaspectratio]{deployment.png}
    \caption{Deployment diagram: ο browser του χρήστη επικοινωνεί με τον
    Vite dev server στο port 5173 (που σερβίρει τα React assets) και με τον
    uvicorn process στο port 8000 (FastAPI). Το uvicorn συνδέεται μέσω
    asyncmy pool στο MySQL 8.0 container, που τρέχει σε Docker network.}
    \label{fig:deployment}
\end{figure}
